\documentclass[11pt]{report}
\usepackage[margin=1in]{geometry}
\usepackage{amsmath, amssymb}
\usepackage{listings}
\usepackage{xcolor}
\usepackage{hyperref}
\usepackage{tcolorbox}
\tcbuselibrary{listings, breakable, skins}

\definecolor{tfblue}{RGB}{0,100,160}
\definecolor{codebg}{RGB}{248,248,248}
\definecolor{commentgreen}{RGB}{0,128,0}
\definecolor{strorange}{RGB}{200,100,0}

\lstdefinestyle{python}{
  language=Python,
  basicstyle=\ttfamily\small,
  keywordstyle=\color{tfblue}\bfseries,
  commentstyle=\color{commentgreen},
  stringstyle=\color{strorange},
  backgroundcolor=\color{codebg},
  showstringspaces=false,
  breaklines=true,
  frame=single,
  framesep=4pt,
}
\lstdefinestyle{shell}{
  basicstyle=\ttfamily\small,
  backgroundcolor=\color{codebg},
  breaklines=true,
  frame=single,
  framesep=4pt,
}

% Box for "new TF/TFP command" callouts.
\newtcolorbox{tfcmd}[1]{
  colback=blue!4, colframe=tfblue, fonttitle=\bfseries,
  title={New command: \texttt{#1}}, breakable
}
% Box for test-result summaries.
\newtcolorbox{testresult}[1]{
  colback=green!4, colframe=green!50!black, fonttitle=\bfseries,
  title={Test result: #1}, breakable
}

\title{JPML\_TF Rebuild Notes}
\author{Haowu Duan}
\date{}

\begin{document}
\maketitle
\tableofcontents

\chapter{Introduction}

This document is a working log of the from-scratch TensorFlow rebuild of
the \texttt{JPMLCOE} codebase under \texttt{code\_clean/jpml\_tf/}. It
serves three purposes:

\begin{enumerate}
  \item Track what was built at each phase, in build order, with the
    exact files involved.
  \item Explain every new TensorFlow or TensorFlow Probability command
    the first time it appears, in plain language. Standard library
    functions (\texttt{json}, \texttt{pathlib}, etc.) are not explained.
  \item Capture the test command and test result for every stage, with
    a short justification of why the result demonstrates the code is
    correct at that point.
\end{enumerate}

The plan being executed is in
\texttt{rebuilt\_plan/00\_overview.md} through
\texttt{08\_migration.md}. Each chapter below corresponds to one phase
in that plan; sections within a chapter correspond to incremental
sub-pieces inside the phase.

\paragraph{Environment.} Python 3.12, TensorFlow 2.16.2,
TensorFlow Probability 0.24.0, Pydantic 2.13.3, pytest 9.0.2,
PyYAML 6.0.3. Package manager: \texttt{uv} 0.11.8. Virtualenv at
\texttt{code/.venv/}. The \texttt{code\_clean/} tree is the new code
root; the legacy \texttt{code/} tree is untouched.

\paragraph{Convention.} Code listings show the file path immediately
above the listing. New TensorFlow commands appear in blue boxes labelled
\textit{New command}. Test results appear in green boxes that include
the exact \texttt{pytest} command, the captured output, and a short
justification.

\chapter{Phase 01 --- Foundation, Config, Results Sink}

Phase~01 builds the package import boundary, the global TensorFlow
startup hook that pins \texttt{float64}, the Pydantic-validated scenario
schema, and the uniform JSON test-result sink. Every later phase loads
through these.

This chapter is built in sub-pieces. Each sub-piece introduces the
smallest possible change, has its own gate test, and is documented
before the next one starts.

\section{Sub-piece 1.1 --- Package import boundary}

\subsection*{Goal}

Establish the package root \texttt{jpml\_tf/} so that
\texttt{import jpml\_tf} succeeds and has no import side effects, in
particular it does not import TensorFlow. The dtype-policy invariant
that later phases depend on requires that no TF code runs before
\texttt{configure\_tf()} is called; that invariant is impossible to
maintain unless the package root stays empty of computational imports.

\subsection*{Files built}

\begin{lstlisting}[style=python, caption={\texttt{code\_clean/jpml\_tf/\_\_init\_\_.py}}]
"""jpml_tf -- TensorFlow rebuild of JPMLCOE.

The package root is intentionally empty: importing ``jpml_tf`` must not
trigger TensorFlow imports, scenario loading, or array creation. Every
entry point (CLI, test, notebook) calls ``jpml_tf.startup.configure_tf()``
explicitly before any computational module is imported.
"""

__version__ = "0.1.0"
\end{lstlisting}

\begin{lstlisting}[style=python, caption={\texttt{code\_clean/jpml\_tf/types.py}}]
"""Shared cross-module types.

Populated incrementally as later phases need them. Phase 01 only owns
the package scaffold and does not yet require any types.
"""
\end{lstlisting}

\begin{lstlisting}[style=python, caption={\texttt{code\_clean/tests/conftest.py}}]
"""pytest configuration shared by all tests under ``code_clean/tests/``.

Adds the package parent (``code_clean/``) to ``sys.path`` so tests can
``import jpml_tf`` without installing the package.
"""

import sys
from pathlib import Path

_PKG_PARENT = Path(__file__).resolve().parent.parent  # code_clean/
if str(_PKG_PARENT) not in sys.path:
    sys.path.insert(0, str(_PKG_PARENT))
\end{lstlisting}

\begin{lstlisting}[style=python, caption={\texttt{code\_clean/tests/test\_01\_imports.py} (excerpt)}]
def test_jpml_tf_imports() -> None:
    import jpml_tf  # noqa: F401


def test_jpml_tf_exposes_version_string() -> None:
    import jpml_tf

    assert isinstance(jpml_tf.__version__, str)
    assert len(jpml_tf.__version__) > 0


def test_jpml_tf_does_not_import_tensorflow() -> None:
    for mod in [m for m in sys.modules
                if m == "tensorflow" or m.startswith("tensorflow.")]:
        del sys.modules[mod]

    import jpml_tf  # noqa: F401

    tf_loaded = [m for m in sys.modules
                 if m == "tensorflow" or m.startswith("tensorflow.")]
    assert tf_loaded == [], (
        f"Importing jpml_tf must not load TensorFlow, but found: {tf_loaded[:5]}"
    )
\end{lstlisting}

\subsection*{New commands introduced}

This sub-piece introduces no TensorFlow or TensorFlow Probability
commands. The only library-level facts in play are standard Python:
\texttt{sys.modules} (the dictionary of every module Python has
imported in the current process) and \texttt{Path(\_\_file\_\_).resolve().parent}
(the absolute directory containing the current source file).

\subsection*{Test command and result}

\begin{lstlisting}[style=shell]
cd code_clean
python -m pytest tests/test_01_imports.py -v
\end{lstlisting}

\begin{testresult}{3 passed, 0.01s}
\begin{lstlisting}[style=shell]
tests/test_01_imports.py::test_jpml_tf_imports                  PASSED [ 33%]
tests/test_01_imports.py::test_jpml_tf_exposes_version_string   PASSED [ 66%]
tests/test_01_imports.py::test_jpml_tf_does_not_import_tensorflow PASSED [100%]
============================== 3 passed in 0.01s ===============================
\end{lstlisting}
\end{testresult}

\subsection*{Why this result is sufficient}

The first two tests are smoke checks: the package imports and exposes
its version string. They would fail if \texttt{\_\_init\_\_.py} had a
syntax error or if \texttt{\_\_version\_\_} were misspelled.

The third test is the load-bearing one. It clears any TensorFlow
modules from \texttt{sys.modules}, imports \texttt{jpml\_tf}, and then
asserts that no TensorFlow module re-appeared. This proves that the
package root has no transitive import path to TensorFlow. As long as
this test passes on every change, the dtype-policy invariant introduced
in the next sub-piece (configure \texttt{float64} \emph{before} TF
creates any tensor) cannot be silently broken by an accidental import.

There are no JSON-saved numerical results yet because the uniform
result sink is built later in this phase (sub-piece~1.3). Once the sink
exists, all phase tests --- including these three --- will route their
results through it.

\section{Sub-piece 1.2 --- TensorFlow startup hook}

\subsection*{Goal}

Provide a single function ``configure\_tf()'' that callers (CLI, test,
notebook) invoke once per process, before any computational module of
\texttt{jpml\_tf} is imported and before any TensorFlow tensor is
created. The function locks the global numeric mode to \texttt{float64},
configures device visibility, and optionally enables deterministic ops.
Phase~01 of the rebuild plan locks \texttt{float64} as the only
critical-path dtype; this sub-piece is where that lock physically lives.

\subsection*{Files built}

\begin{lstlisting}[style=python, caption={\texttt{code\_clean/jpml\_tf/startup.py}}]
"""Process-level TensorFlow configuration."""

from __future__ import annotations

import os

_CONFIGURED = False


def configure_tf(
    *,
    force_cpu: bool = False,
    allow_gpu_growth: bool = True,
    deterministic_ops: bool = False,
    legacy_keras: bool = True,
) -> None:
    """Lock global TF numeric and device policy to rebuild defaults."""
    global _CONFIGURED

    os.environ["TF_USE_LEGACY_KERAS"] = "1" if legacy_keras else "0"

    import tensorflow as tf

    if force_cpu:
        tf.config.set_visible_devices([], "GPU")
    elif allow_gpu_growth:
        for gpu in tf.config.list_physical_devices("GPU"):
            try:
                tf.config.experimental.set_memory_growth(gpu, True)
            except RuntimeError:
                pass

    if deterministic_ops:
        tf.config.experimental.enable_op_determinism()

    tf.keras.backend.set_floatx("float64")
    tf.keras.mixed_precision.set_global_policy("float64")

    _CONFIGURED = True


def is_configured() -> bool:
    return _CONFIGURED
\end{lstlisting}

\begin{lstlisting}[style=python, caption={\texttt{code\_clean/tests/test\_01\_startup.py}}]
import os


def test_configure_tf_sets_float64() -> None:
    from jpml_tf.startup import configure_tf
    configure_tf()

    import tensorflow as tf
    assert tf.keras.backend.floatx() == "float64"
    assert tf.keras.mixed_precision.global_policy().compute_dtype == "float64"


def test_configure_tf_is_idempotent() -> None:
    from jpml_tf.startup import configure_tf, is_configured
    configure_tf()
    configure_tf()
    assert is_configured()

    import tensorflow as tf
    assert tf.keras.backend.floatx() == "float64"


def test_configure_tf_sets_legacy_keras_env_var() -> None:
    from jpml_tf.startup import configure_tf
    configure_tf(legacy_keras=True)
    assert os.environ.get("TF_USE_LEGACY_KERAS") == "1"
\end{lstlisting}

\subsection*{New commands introduced}

\begin{tfcmd}{tf.keras.backend.set\_floatx, tf.keras.backend.floatx}
TensorFlow's Keras submodule has a process-wide default float dtype
called \emph{floatx}. It controls the dtype of any Keras layer, Keras
optimizer state, or anything that asks Keras ``what's your default
float?'' \texttt{set\_floatx("float64")} flips that default;
\texttt{floatx()} reads the current value. Default at TF startup is
\texttt{"float32"}.

\textbf{Caveat to remember.} \emph{floatx} only governs Keras-aware code.
Plain \texttt{tf.constant(1.0)} still gives \texttt{float32} regardless
of \texttt{floatx}, because \texttt{tf.constant} infers dtype from the
literal's Python type. The rebuild's discipline is to always pass
explicit \texttt{dtype=tf.float64} when constructing tensors; setting
\emph{floatx} catches the Keras-side defaults that explicit-dtype
discipline cannot reach.
\end{tfcmd}

\begin{tfcmd}{tf.keras.mixed\_precision.set\_global\_policy, .global\_policy}
The mixed-precision API in Keras lets you say ``compute in this dtype,
store variables in that dtype.'' Setting the \emph{global policy} to
\texttt{"float64"} forces Keras to use float64 for both compute and
variables. It is a stronger lock than \emph{floatx}: any Keras layer
that asks ``what compute dtype should I use?'' gets ``float64'' from
the policy, even if the layer was constructed with no explicit
\texttt{dtype} argument.

\textbf{Why both \emph{floatx} and the policy.} They control overlapping
but not identical surfaces. \emph{floatx} is the older API, used by
some optimizers and helpers. The policy API is the newer, broader
mechanism. Setting both belt-and-braces is cheap and avoids ``it works
in some versions of Keras but not others'' surprises.
\end{tfcmd}

\begin{tfcmd}{tf.config.list\_physical\_devices, tf.config.set\_visible\_devices}
\texttt{list\_physical\_devices("GPU")} returns the list of GPUs TF can
see on this machine, as a list of \texttt{PhysicalDevice} objects.

\texttt{set\_visible\_devices([], "GPU")} hides all GPUs from TF for
the rest of the process: every subsequent op runs on CPU. This is the
clean way to force CPU-only mode --- much more reliable than setting
\texttt{CUDA\_VISIBLE\_DEVICES} after TF has already initialised.

\textbf{Order matters.} Both calls must happen before TF places any op
on a device. After the first computation runs, TF locks the device
visibility for the process; later calls error out.
\end{tfcmd}

\begin{tfcmd}{tf.config.experimental.set\_memory\_growth}
By default, TF eagerly allocates almost all of a GPU's memory the
moment it touches that GPU. \texttt{set\_memory\_growth(gpu, True)}
asks TF to grow its memory pool incrementally instead, which lets
multiple TF processes share a GPU and prevents spurious OOM errors
from one greedy process. The trade-off is a small fragmentation cost.
The \texttt{experimental} prefix means the API name might change in a
future TF version, but it has been in this form since TF~2.0.
\end{tfcmd}

\begin{tfcmd}{tf.config.experimental.enable\_op\_determinism}
Some TF ops (e.g. cuDNN convolutions, segment reductions, parallel
reductions) are intentionally non-deterministic so they can use faster
parallel implementations. \texttt{enable\_op\_determinism()} forces TF
to use deterministic implementations everywhere. The result: the same
input plus the same seed always produces the same output, at the cost
of slower training. Useful for debugging and for paper-quality
reproducibility runs; disabled by default in the rebuild because most
training is faster without it.
\end{tfcmd}

\begin{tfcmd}{TF\_USE\_LEGACY\_KERAS environment variable}
Not a TF command, but the only knob that controls which Keras backend
TensorFlow uses. TF~2.16+ ships Keras~3 by default; setting
\texttt{TF\_USE\_LEGACY\_KERAS=1} switches it to the older Keras~2.x
backend (\texttt{tf-keras}, installed separately). The rebuild defaults
to legacy Keras for compatibility with the existing TFP HMC kernels,
which still target the Keras~2.x backend on the office RTX~3090 box.
TF reads this variable on its first import; setting it later has no
effect.
\end{tfcmd}

\subsection*{Test command and result}

\begin{lstlisting}[style=shell]
cd code_clean
python -m pytest tests/ -v
\end{lstlisting}

\begin{testresult}{6 passed, 2.56s}
\begin{lstlisting}[style=shell]
tests/test_01_imports.py::test_jpml_tf_imports                       PASSED
tests/test_01_imports.py::test_jpml_tf_exposes_version_string        PASSED
tests/test_01_imports.py::test_jpml_tf_does_not_import_tensorflow    PASSED
tests/test_01_startup.py::test_configure_tf_sets_float64             PASSED
tests/test_01_startup.py::test_configure_tf_is_idempotent            PASSED
tests/test_01_startup.py::test_configure_tf_sets_legacy_keras_env_var PASSED
================== 6 passed, 2 warnings in 2.56s =====================
\end{lstlisting}
\end{testresult}

\subsection*{Why this result is sufficient}

The first new test asserts the two complementary dtype locks took
effect: \texttt{floatx()} is \texttt{"float64"} \emph{and} the
mixed-precision compute dtype is \texttt{"float64"}. Either one alone
would leave gaps that the other closes; both together cover every
Keras code path the rebuild touches.

The second test asserts idempotence. Real entry points may call
\texttt{configure\_tf()} more than once --- a script that imports a
test fixture which itself calls it, for example. Idempotence guarantees
this is safe: nothing throws, and the dtype stays locked.

The third test asserts that \texttt{TF\_USE\_LEGACY\_KERAS=1} actually
landed in \texttt{os.environ}. This is a surface check: the env var is
the only public signal that controls Keras-2 versus Keras-3 mode, and
later phases assume Keras~2 semantics for layer dtype propagation.

The test result is currently captured only as pytest stdout in this
document. Sub-piece~1.3 builds the uniform JSON result sink; once it
exists, all six tests above will be re-run with their assertions
recorded as JSON records keyed by test name.

\subsection*{State of \texttt{code\_clean/} after this sub-piece}

\begin{lstlisting}[style=shell]
code_clean/
  jpml_tf/
    __init__.py
    startup.py               (new)
    types.py
  tests/
    __init__.py
    conftest.py
    test_01_imports.py
    test_01_startup.py       (new)
  notes/
    notes.tex                (this document)
    notes.pdf                (compiled output)
\end{lstlisting}

\end{document}
